%% units.tex -- condition (U) as a plain single-word hole problem: the source
%% automaton, the unit reader with the rejected-path property, and the decision
%% of (U) by the hole search of Zhao26 on the unit reader. The smaller reader
%% and the antecedent rejection clock, used only for the complexity bounds, are
%% in the unmechanized part (complexity.tex, sec:un-smaller). Proofs:
%% appendix-summaries.tex (app:sm).

\section{Condition (U)}\label{sec:un}

From now on, $\varphi$ is a system over the visible variables $X$ and $\psi$ a
system over $X\cup Z$, as in \cref{sec:ch}. Condition \CondU\ of
\cref{def:ch-conditions} asks that $\Reach_\psi(a)$ contain no forbidden root of
$\psi$ for every visible atom $a$ that is live in $\varphi$. In the terms of
\cref{def:sm-deadu}, this reads as follows.

\begin{lemma}[Condition (U)]\label{lem:un-dead}
Condition \CondU\ holds iff no $p\in X^\pm$ and no word $w$ satisfy
$\DeadU_\psi(p,w)$ while $(p,w)$ is live in $\varphi$.
\end{lemma}

The words $w$ with $\DeadU_\psi(p,w)$ form a regular language
(\cref{lem:sm-dead-states}), but the words $w$ with $(p,w)$ live in $\varphi$
need not (\cref{ex:sm-arrow}), so \CondU\ is not an inclusion of regular
languages. When $\varphi$ is satisfiable, however, \cref{thm:sm-dead}
identifies liveness with the absence of $\DeadT_\varphi$, which is regular,
and of self events, which are comparisons of a single word whose two cuts are
tested by a pair of automaton states. A failure of \CondU\ is therefore a hole of a reader in the sense of
\citet{Zhao26} (\cref{def:bg-reader,def:bg-hole}), and the reader has the
rejected-path property (\cref{def:bg-rpp}). The hole search of that paper
(\cref{cor:bg-hole-decision}) then decides \CondU.

\subsection{The Source Automaton}\label{sec:un-automaton}

One automaton reads the word of a visible seed for both systems. The two-word
instances of \cref{sec:ex} use it again, on the seed word.

\begin{definition}[Source automaton]\label{def:un-source}
Let $p\in X^\pm$. The \emph{source automaton} $\Bsrc_p$ is the deterministic
automaton over $\Sigma$ whose states are the triples $\zeta=(P,b,P')$ with
$P\subseteq X^\pm$, $b\in\{0,1\}$, and $P'\subseteq(X\cup Z)^\pm$, with start
state $\bigl(\upc\varphi p,\ [\overline{\upc\varphi p}\cap\Fup\varphi\ne
\varnothing],\ \upc\psi p\bigr)$ and transitions
\[
 \delta\bigl((P,b,P'),i\bigr)=\bigl(\pstep i(P),\
 b\lor[\overline{\pstep i(P)}\cap\Fup\varphi\ne\varnothing],\
 \pstep i(P')\bigr),
\]
where the first component uses the transitions of the upper automata of
$\varphi$ and the third those of $\psi$, and $[\,\cdot\,]$ is $1$ for a true
condition and $0$ otherwise. For a state $\zeta=(P,b,P')$ put
\[
 \mathrm{dt}(\zeta)\iff b=1\ \text{or}\ P\cap\DeadT^\eps_\varphi\ne\varnothing,
 \qquad
 \mathrm{du}(\zeta)\iff P'\cap\DeadU^\eps_\psi\ne\varnothing,
\]
and let the \emph{self guard} $\GuardS$ consist of the pairs of states
$((P_1,b_1,P'_1),(P_2,b_2,P'_2))$ with $P_1\cap\overline{P_2}\ne\varnothing$.
\end{definition}

\begin{lemma}[Runs of the source automaton]\label{lem:un-source}
Let $p\in X^\pm$, let $w$ be a word, and let $\zeta_j$ be the state of
$\Bsrc_p$ after $\wpre wj$. Then:
\begin{enumerate}[leftmargin=*]
\item $\zeta_j=\bigl(P^\varphi_{p,j}(w),\ [(\bar p,\wpre wj)\in\mu_\varphi],\
  P^\psi_{p,j}(w)\bigr)$;
\item $\mathrm{dt}(\zeta_j)$ iff $\DeadT_\varphi(p,\wpre wj)$, and
  $\mathrm{du}(\zeta_j)$ iff $\DeadU_\psi(p,\wpre wj)$; if
  $\mathrm{dt}(\zeta_j)$ fails, it fails at every later cut, and if
  $\mathrm{du}(\zeta_j)$ holds, it holds at every earlier cut;
\item for $i<j\le\len w$, the pair $(i,j)$ is a self event of $(p,w)$ in
  $\varphi$ iff it is a comparison of $w$ and $(\zeta_i,\zeta_j)\in\GuardS$.
\end{enumerate}
The automaton $\Bsrc_p$ has $N_{\mathrm U}=2^{4\len X+2\len Z+1}$ states.
\end{lemma}

Item (1) holds by induction on $j$, using \cref{lem:sm-dead-states}(2) for
the bit; item (2) is \cref{lem:sm-dead-states}(1,3) together with
\cref{lem:sm-roots}(4) and \cref{lem:sm-dead-states}(4); and item (3) is
\cref{def:sm-dead}. The sets $\Fup\varphi$, $\DeadT^\eps_\varphi$, and
$\DeadU^\eps_\psi$ are computable (\cref{lem:sm-mu,lem:sm-dead-states}), so
$\Bsrc_p$, $\mathrm{dt}$, $\mathrm{du}$, and $\GuardS$ are computable.

\subsection{The Unit Reader}\label{sec:un-reader}

\begin{definition}[Unit reader]\label{def:un-reader}
For $p\in X^\pm$, the \emph{unit reader} is the reader
$\RU(p)=(\Bsrc_p,\GuardS,\Rej_{\mathrm U})$ with the source automaton, the
self guard as its admission relation, and the rejected states
\[
 \Rej_{\mathrm U}=\{\zeta:\ \mathrm{dt}(\zeta)\text{ fails and }
 \mathrm{du}(\zeta)\text{ holds}\}.
\]
\end{definition}

\begin{lemma}[Holes of the unit reader]\label{lem:un-holes}
Let $\varphi$ be satisfiable and $p\in X^\pm$. A word $w$ is a hole of
$\RU(p)$ iff $(p,w)$ is live in $\varphi$ and $\DeadU_\psi(p,w)$.
Consequently, \CondU\ holds iff none of the $2\len X$ unit readers has a hole.
\end{lemma}

\begin{proof}
By \cref{lem:un-source}, $w$ is rejected iff $\DeadT_\varphi(p,w)$ fails and
$\DeadU_\psi(p,w)$ holds, and a comparison of $w$ is admitted iff it is a self
event of $(p,w)$ in $\varphi$. So $w$ is a hole iff $\DeadU_\psi(p,w)$ holds
and $(p,w)$ is neither dead by support nor has a self event in $\varphi$, that
is, by \cref{thm:sm-dead}, iff $\DeadU_\psi(p,w)$ holds and $(p,w)$ is live in
$\varphi$. The second claim follows by \cref{lem:un-dead}.
\end{proof}

\begin{lemma}[Rejected paths]\label{lem:un-rejected-path}
Every unit reader $\RU(p)$ has the rejected-path property: it satisfies
betweenness and the admission cone. No satisfiability assumption is needed.
\end{lemma}

\begin{proof}
\emph{Betweenness.} Let $\wpre wa$ and $\wpre wb$ be rejected, with $a<k<b$.
Since $\mathrm{dt}$ fails at $a$, it fails at $k$, and since $\mathrm{du}$
holds at $b$, it holds at $k$ (\cref{lem:un-source}(2)). So $\wpre wk$ is
rejected.

\emph{Admission cone.} Let $w$ be rejected, with first rejected cut $J$, and
let $(i,j)$ be an admitted comparison of $w$. The guard
$P^\varphi_{p,i}(w)\cap\overline{P^\varphi_{p,j}(w)}\ne\varnothing$ requires
$P^\varphi_{p,j}(w)\ne\varnothing$. At the cut $J$, $\mathrm{dt}$ fails, so
$\DeadT_\varphi(p,\wpre wJ)$ fails, and in particular
$\DeadU_\varphi(p,\wpre wJ)$ fails. The pop bound (\cref{lem:sm-roots}(2)) gives
$P^\varphi_{p,j'}(w)=\varnothing$ for every $j'>J$. Hence $j\le J$.
\end{proof}

The proof gives more than the admission cone: no admission, comparison or not, has its
second cut after $J$.

\begin{example}[Shift closure: condition (U)]\label{ex:un-running}
Let $\psi=\{x\le z,\ z\le y,\ z\le f(z,z)\}$ as in \cref{ex:sm-running}, with
$f$ binary and covariant, and let $\varphi_1=\{x\le y\}$. Over $\BFT$,
condition \CondU\ fails. Indeed, $\DeadU^\eps_\psi=\{x^0,z^0\}$ and
$z^0\in P^\psi_{x^0}(w)$ for every $w$, so $\DeadU_\psi(x^0,w)$ holds for every
word $w$. On the other side, $\varphi_1$ has neither facts nor
forbidden roots, so $\DeadT_{\varphi_1}$ never holds, and it has no upper
constructor literal, so $P^{\varphi_1}_{x^0}(w)=\varnothing$ for $w\ne\eps$ and
no comparison is admitted. Every word is a hole of $\RU(x^0)$, the empty word
included. Its countermodel is the least model $\Cl_{\varphi_1}(x^0,\eps)$, the
solution $x=y=\Top$: no $z$ satisfies $\Top\le z\le f(z,z)$.

Over $\BF$, the variables cannot take the value $\Top$. Within $\BFT$, this
can be expressed by a variable $g$ with $y\le g$ and $g\le f(g,g)$, which
allow exactly the trees $y$ without leaves $\Top$; this is the gadget of
\cref{lem:tc-gadget}. The reduction of \cref{thm:tc-gadget} puts such a gadget
on every variable, but for this example one on $y$ suffices: since $x\le y$, it also
keeps $x$ free of $\Top$, and $z\le f(z,z)$ already keeps $z$ free of $\Top$.
So let $\varphi'_1=\{x\le y,\ y\le g,\ g\le f(g,g)\}$, with $g$ visible. Its only
forbidden root is $(g^0,\eps)$, and $x^0\to y^0\to g^0$ by copies, after which
$g^0$ pops every word. So $\DeadT_{\varphi'_1}(x^0,w)$ holds for every $w$, and
$\RU(x^0)$ rejects no word. For the other signed variables
$p\in\{x^1,y^0,y^1,g^0,g^1\}$, the set $P^\psi_p(w)$ never meets
$\DeadU^\eps_\psi$, so $\RU(p)$ rejects no word either. Thus \CondU\ holds for
$\varphi'_1$; \cref{ex:ex-running} shows that the entailment fails by
condition \CondII.
\end{example}

\begin{example}[Admitted self events are needed]\label{ex:un-arrow}
Let $\varphi=\{x\le x\to x\}$ as in \cref{ex:sm-arrow}, and let
$\psi=\{x\le a\to x,\ c\to y\le a\}$ with $a,c$ hidden. The antecedent is that
of \cref{ex:ch-arrow}, but the consequent is new. For $x\le z\to z$,
\CondU\ holds because its only root-dead atom, $(x^0,\eps)$, is a forbidden
root of $\varphi$ as well; this consequent is chosen so that \CondU\ holds only
because of admitted self events. The signed system of
$\psi$ has the upper constructor literals $x^0\le f(a^1,x^0)$ and
$a^1\le f(c^0,y^1)$ and the lower ones $f(a^0,x^1)\le x^1$ and
$f(c^1,y^0)\le a^0$; its closure is the identity, and
$\DeadU^\eps_\psi=\{x^0,a^1\}$. Reading from $x^0$,
\[
 P^\psi_{x^0}(1^k)=\{x^0\},\quad P^\psi_{x^0}(1^k0)=\{a^1\},\quad
 P^\psi_{x^0}(1^k00)=\{c^0\},\quad P^\psi_{x^0}(1^k01)=\{y^1\},
\]
and the empty set after longer words. So $\DeadU_\psi(x^0,w)$ holds iff
$w\in1^*\cup1^*0$, and $\RU(x^0)$ rejects exactly the words $1^k0$, which are
not dead by support in $\varphi$. Each of them has the admitted comparison
$(0,k+1)$: $P^\varphi_{x^0}(\eps)=\{x^0\}$ meets
$\overline{P^\varphi_{x^0}(1^k0)}=\{x^0\}$. So $\RU(x^0)$ has no hole. The other
unit readers reject no word, because $x^1$, $y^0$, and $y^1$ have no root-dead
atoms in $\psi$. Hence \CondU\ holds. Without the admitted self events, the
word $0$ would be a hole: $(x^0,0)$ is dead in $\varphi$ only through its self event
(\cref{ex:sm-arrow}).
\end{example}

\Cref{fig:un-reader} draws this reader. The admission relation acts as a
rescue: a rejected word, a candidate failure of \CondU, is discarded as soon
as two of its cuts form a comparison with an admitted pair of states, because
the seed is then dead in $\varphi$. A small change of the consequent, also
shown there, leaves a rejected word that no comparison rescues, a hole.

\begin{figure}[tbp]
\centering
\begin{tikzpicture}[>=Stealth,font=\small,x=1cm,y=1cm,
    every node/.style={inner sep=1.5pt},
    st/.style={state,minimum size=6.5mm,inner sep=0pt},
    lb/.style={font=\scriptsize,inner sep=1pt},
    cell/.style={draw,minimum width=6mm,minimum height=5mm,inner sep=0pt},
    adm/.style={thick},
    nadm/.style={densely dashed}]
  % the reduced unit reader R_U(x^0)
  \node[st,initial,initial text={}] (A) at (0,0) {$A$};
  \node[st,accepting] (B) at (2.2,0) {$B$};
  \node[st,accepting,dashed] (D) at (4.2,0) {$D$};
  \node[st] (S) at (3.2,-1.5) {$\bot$};
  \path[->] (A) edge[loop above] node[lb,above]{$1$} (A)
            (A) edge node[lb,above]{$0$} (B)
            (B) edge node[lb,above]{$1$} (D)
            (B) edge node[lb,left]{$0$} (S)
            (D) edge node[lb,right]{$0,1$} (S)
            (S) edge[loop left] node[lb,left]{$0,1$} (S);
  \draw[<->,adm,densely dotted] (A) to[bend right=40] node[lb,below]{admitted} (B);
  \node[lb,anchor=north west,align=left] at (-0.45,-1.25)
    {$A=(\{x^0\},0,\{x^0\})$\\ $B=(\{x^1\},0,\{a^1\})$\\ $D=(\varnothing,0,\{y^1\})$};
  % the word 110
  \begin{scope}[xshift=6.6cm]
    \node[font=\footnotesize,anchor=west] at (-0.3,1.45) {$110$: rejected, but not a hole};
    \foreach \l [count=\i from 0] in {1,1,0}
      \node[cell] at (0.6*\i+0.3,0) {$\l$};
    \foreach \q [count=\i from 0] in {A,A,A,B}
      \node[lb,anchor=north] at (0.6*\i,-0.3) {$\q$};
    \draw[adm] (0,0.25) .. controls (0.3,1.15) and (1.5,1.15) .. (1.8,0.25);
    \draw[adm] (0.6,0.25) .. controls (0.85,0.85) and (1.55,0.85) .. (1.8,0.25);
    \draw[adm] (1.2,0.25) .. controls (1.35,0.6) and (1.65,0.6) .. (1.8,0.25);
    \node[lb,anchor=west,align=left] at (2.0,0.25) {$(0,3)$, $(1,3)$, $(2,3)$:\\ pairs $(A,B)$, admitted};
  \end{scope}
  % the word 01, variant
  \begin{scope}[xshift=6.6cm,yshift=-2.15cm]
    \node[font=\footnotesize,anchor=west] at (-0.3,1.2) {$01$ in the variant: a hole};
    \foreach \l [count=\i from 0] in {0,1}
      \node[cell] at (0.6*\i+0.3,0) {$\l$};
    \foreach \q [count=\i from 0] in {A,B,D}
      \node[lb,anchor=north] at (0.6*\i,-0.3) {$\q$};
    \draw[nadm] (0,0.25) .. controls (0.25,0.95) and (0.95,0.95) .. (1.2,0.25);
    \draw[nadm] (0.6,0.25) .. controls (0.75,0.6) and (1.05,0.6) .. (1.2,0.25);
    \node[lb,anchor=west,align=left] at (1.4,0.25) {$(0,2)$, $(1,2)$: pairs $(A,D)$, $(B,D)$,\\ not admitted; $(0,1)$ is no comparison};
  \end{scope}
\end{tikzpicture}
\caption{The unit reader $\RU(x^0)$ of \cref{ex:un-arrow}, after merging the
states that cannot occur on a rejected word into a sink $\bot$
(\cref{lem:un-reduced}). Rejected states are doubly circled; the only admitted
pairs are $(A,B)$ and $(B,A)$, the self guard of the antecedent $x\le x\to x$.
Every rejected word $1^k0$, such as $110$, ends in $B$ and has the terminal
comparisons $(i,k+1)$ with the admitted pair $(A,B)$: it is rescued, because
$(x^0,1^k0)$ is dead in $\varphi$ by a self event. In the example itself the
state $D$, dashed, is not rejected and would be merged into $\bot$. In a
variant in which $y$ is hidden and $\psi$ also contains $d\to e\le y$, with
$d,e$ hidden, the atom $(y^1,\eps)$ is a forbidden root, $D$ is rejected, and
the word $01$ is a hole: none of its comparisons is admitted, so $(x^0,01)$ is
live in $\varphi$ and root-dead in $\psi$, and \CondU\ fails.}
\Description{Left: an automaton with states A, B, D and a sink bottom. A is the
start state with a loop on 1; 0 leads from A to B, 1 from B to D, 0 from B to
the sink, and both letters from D to the sink, which loops on both letters. B
is doubly circled, D doubly circled and dashed. A dotted double arrow labeled
admitted joins A and B. A legend lists the triples of the states A, B and D. Right: the
word 110 with the states A, A, A, B at its cuts and three solid arcs from the
cuts 0, 1 and 2 to the cut 3, labeled admitted; below, the word 01 with the states
A, B, D and two dashed arcs from the cuts 0 and 1 to the cut 2, labeled not
admitted.}
\label{fig:un-reader}
\end{figure}


\subsection{Deciding Condition (U)}\label{sec:un-decide}

\Citet{Zhao26} decides the hole problem for every reader with the
rejected-path property (\cref{cor:bg-hole-decision}): a reader with $N$ states
over a finite alphabet that has this property has a hole iff it has a hole of
length at most $\holeBound(N,d_N)$, where $d_N=2N(4N-1)-1$. The proof selects a
hole with rejected tail at most $4N-2$ (\cref{thm:bg-brt}), bounds its depth
(\cref{thm:bg-depth}), and shortens it (\cref{thm:bg-search}). Applied to the
unit readers, with their $N_{\mathrm U}$ states, this decides \CondU.

\begin{theorem}[Condition (U) is decidable]\label{thm:un-decidable}
Let $\varphi$ be satisfiable, and let $N_{\mathrm U}=2^{4\len X+2\len Z+1}$ be
the number of states of the source automaton. Condition \CondU\ fails iff
there are $p\in X^\pm$ and a word $w$ with
$\len w\le\holeBound(N_{\mathrm U},d_{N_{\mathrm U}})$ such that $(p,w)$ is
live in $\varphi$ and $\DeadU_\psi(p,w)$. Hence \CondU\ is decidable. If
$X\ne\varnothing$, then, with $v=\len X$, $s=\len X+\len Z$ and
$E_3(x)=2^{2^{2^x}}$,
\[
 \holeBound(N_{\mathrm U},d_{N_{\mathrm U}})\le E_3(10v+10s+17)\le E_3(20s+17).
\]
The bound depends neither on the arity nor on the variance.
\end{theorem}

\begin{proof}
By \cref{lem:un-holes}, \CondU\ fails iff some unit reader $\RU(p)$ has a
hole. Each $\RU(p)$ has $N_{\mathrm U}$ states, the finite alphabet $\Sigma$,
and the rejected-path property (\cref{lem:un-rejected-path}), so by
\cref{cor:bg-hole-decision} it has a hole iff it has one of length at most
$\holeBound(N_{\mathrm U},d_{N_{\mathrm U}})$, and by \cref{lem:un-holes} such
a hole is a word as stated. If $X$ is empty, there is no unit reader and
\CondU\ holds. For the decision, enumerate the $2\len X$ signed variables and
the finitely many words up to this length, and test each pair by
\cref{lem:un-source}: the states of $\Bsrc_p$ along $w$ are computable, and so
are rejection and the admission of each of the finitely many comparisons of
$w$.

For the estimate, let $v\ge1$. Then $N_{\mathrm U}=2^{2v+2s+1}\ge8$, so
$N_{\mathrm U}+1\le\frac98N_{\mathrm U}$, and $(\frac98)^5<2$. By
\cref{cor:bg-hole-decision}, $\holeBound(N_{\mathrm U},d_{N_{\mathrm U}})\le
2^{2^{1280(N_{\mathrm U}+1)^5}}$, and
$1280(N_{\mathrm U}+1)^5<2^{12}N_{\mathrm U}^5=2^{10v+10s+17}$; finally
$v\le s$.
\end{proof}

The bound of \cref{thm:un-decidable} is not the smallest one. A smaller reader
gives a sharper bound (\cref{prop:un-sharper}); it serves only the complexity
analysis and is therefore placed in the part that is not mechanized
(\cref{sec:un-smaller}).

Informally, a different route would reduce \CondU\ to plain entailment: the
words $w$ with $\DeadU_\psi(p,w)$ form a regular language $W_p$, and \CondU\
asks that in every solution of $\varphi$ the leaf cone of $p$ avoid $W_p$, a
comparison with a regular ground tree determined by $W_p$. We do not pursue
this route; the unit reader gives an explicit witness bound without additional
variables. \Cref{sec:ex} turns to condition \CondII.
